%======================================================================
\section{Theoretical Analysis}
\label{sec:theory}
%======================================================================

\begin{theorem}[Existence of Stackelberg Equilibrium]
\label{thm:existence}
The AEIC-Stackelberg game $\mathcal{G}$ admits a Stackelberg equilibrium $(\theta^*, s^*)$.
\end{theorem}

\begin{proof}
The follower strategy space $\{s^{\text{ind}}, s^{\text{syc}}\}$ is finite, so an optimal response exists for any $\theta$. The parameter space $\Theta$ is compact (bounded). The leader's objective $\text{Accuracy}(\theta) + \lambda \cdot \text{IC}(\theta)$ is bounded on $\Theta$. By the Weierstrass extreme value theorem, the maximum is attained at some $\theta^* \in \Theta$. Together with the followers' best responses, this constitutes a Stackelberg equilibrium.
\end{proof}

\begin{theorem}[$\epsilon$-Incentive Compatibility]
\label{thm:ic}
For any $\epsilon \in (0,1)$, setting $\lambda \geq \frac{1}{1-\epsilon}$ in Eq.~\eqref{eq:leader_obj} guarantees $\text{IC}(\theta^*) \geq \epsilon$.
\end{theorem}

\begin{proof}
We first establish feasibility: there exists $\theta_0 \in \Theta$ with $\text{IC}(\theta_0) = 1$ (set IC/EQ exponents large, peer pressure $\gamma = 0$, making weights strictly monotone in quality).

By optimality of $\theta^*$:
\begin{equation}
    \text{Acc}(\theta^*) + \lambda \cdot \text{IC}(\theta^*) \geq \text{Acc}(\theta_0) + \lambda \cdot 1
\end{equation}

Rearranging with $\text{Acc}(\theta_0) - \text{Acc}(\theta^*) \geq -1$:
\begin{equation}
    \text{IC}(\theta^*) \geq 1 - \frac{1}{\lambda}
\end{equation}

Setting $\lambda \geq \frac{1}{1-\epsilon}$ yields $\text{IC}(\theta^*) \geq \epsilon$.
\end{proof}

\begin{theorem}[Individual Rationality and Welfare]
\label{thm:ir}
(a) Under $\theta^*$, all agents achieve non-negative utility when $R \geq C_{\text{ind}} / \min_i w_i$. \\
(b) $\text{Accuracy}(\theta^*) \geq \text{Accuracy}(\theta_{\text{uniform}})$ for sufficiently small $\lambda$.
\end{theorem}

\begin{proof}
(a) Since $w_i \in [0,1]$ and $C(s_i) \leq C_{\text{ind}}$, choosing $R \geq C_{\text{ind}} / \min_i w_i$ ensures $U_i \geq 0$.

(b) Uniform voting corresponds to $\theta_{\text{unif}} \in \Theta$. By optimality: $\text{Acc}(\theta^*) + \lambda \cdot \text{IC}(\theta^*) \geq \text{Acc}(\theta_{\text{unif}}) + \lambda \cdot \text{IC}(\theta_{\text{unif}})$. For $\lambda \to 0$, this directly yields $\text{Acc}(\theta^*) \geq \text{Acc}(\theta_{\text{unif}})$.
\end{proof}

\paragraph{Convergence of Distributed Consensus.}
Under the Stackelberg dynamics with peer pressure $\gamma > 0$ and communication topology $W$, the consensus energy $E(t) = \sum_{(i,j) \in \mathcal{E}} (\theta_i(t) - \theta_j(t))^2$ satisfies:
\begin{equation}
    E(t) \leq E(0) \cdot \exp\!\left(-\frac{2\gamma t}{n-1}\right)
\end{equation}
guaranteeing exponential convergence to consensus at a rate determined by the algebraic connectivity of $W$ and peer pressure coefficient $\gamma$.
